% =====================================================================
% chapters/07_conclusion.tex
% Chapter 7 — Conclusion and Future Work
% Project: Vision-Based Running Technique Classification from In-the-Wild
%          Internet Videos using Markerless Pose Estimation and Deep Learning
%
% Headline metrics quoted below are the controlled results established in
% Chapter 5; they are not re-derived here. No claim of clinical,
% professional, or real-world deployment readiness is made.
% =====================================================================

\chapter{Conclusion and Future Work}
\label{ch:conclusion}

% ---------------------------------------------------------------------
\section{Conclusion}
\label{sec:concl-summary}

This thesis addressed the problem of classifying running technique as correct or
false from \emph{in-the-wild} internet videos using markerless pose estimation
and deep learning. The work developed a complete pipeline consisting of
markerless pose estimation, landmark filtering and normalisation, motion-signal
extraction, running-cycle detection, eight-phase sampling, landmark-based or
image-based cycle representation, deep-learning classification, and evaluation
under a fixed split by video identifier.

The experimental design separated two questions that are often mixed together:
whether the pose representation is learnable when cycle boundaries are reliable,
and how much performance is lost when cycle boundaries must be detected
automatically. The handmade landmark pipeline, using controlled cycle
segmentation, reached a test accuracy of 0.9107 and a macro-F1 of 0.8991. This
result is interpreted as an upper bound, not as a deployment baseline. The
autocorrelation-based landmark pipeline, using automatic cycles detected from
the \texttt{contact\_diff} signal, reached a test accuracy of 0.5455 and a
macro-F1 of 0.5417 under the reported baseline setting. Because the two
pipelines use the same general pose-representation family but differ in cycle
source, the gap between them localises the main difficulty of the task in the
automatic cycle-detection stage. This interpretation is supported by an
additional fair-gap robustness check: under the same MLP classifier and
5-fold video-level cross-validation, handmade cycles achieved
$0.860 \pm 0.016$ macro-F1, whereas autocorr cycles achieved
$0.615 \pm 0.081$, giving a fair gap of 0.245 macro-F1. A 100-permutation test
on the autocorr MLP setting did not reach statistical significance
($p = 0.2475$), suggesting that the autocorr-derived features are not yet
robustly separable under that simple classifier.

The ablation studies and extended analyses refine this conclusion. The model
ablation showed that architecture choice matters less than segmentation quality:
the CNN-BiLSTM is strongest on clean handmade cycles, while the simpler MLP is
more stable as the reported automatic baseline after validation-based threshold
selection. The feature ablation showed that compact selected-landmark
representations are sufficient and competitive with larger landmark sets. The
frames-per-cycle ablation showed that eight-phase sampling is an efficient
choice, since increasing the number of sampled frames did not improve macro-F1.

The signal ablation showed a trade-off between detection coverage and
semantically meaningful foot-strike-to-foot-strike boundaries; \texttt{contact\_diff}
was therefore retained as the main baseline signal even though \texttt{hip\_y}
achieved the highest classification macro-F1 in the signal ablation. The
controlled landmark-versus-image comparison showed that the RGB-crop
MobileNetV2--LSTM pipeline did not reproduce the historical image-based
development result and was not stable enough to replace the landmark baseline.

Taken together, the results support a diagnostic conclusion. Markerless pose
sequences can encode correct-versus-false running-technique differences when
cycles are clean, but the current automatic pipeline is limited by pose quality,
cycle detection, and boundary alignment. The fixed-split result and the
cross-validation robustness check point to the same conclusion, while the
permutation test shows that the autocorr representation remains statistically
fragile under a simple MLP. The thesis therefore does not claim
real-world deployment readiness. Its contribution is to identify where the
automatic in-the-wild pipeline succeeds, where it fails, and which components
should be improved before any practical feedback system is attempted.

% ---------------------------------------------------------------------
\section{Answers to Research Questions}
\label{sec:concl-rqs}

The seven research questions introduced in Section~\ref{sec:intro-rqs} are
answered below using the controlled results from Chapter~\ref{ch:results}.

\paragraph{RQ1 --- Can pose-based classifiers distinguish correct from false
technique under reliable segmentation?}
Yes. Under reliable manually controlled cycle segmentation, the handmade
landmark upper bound reached a test accuracy of 0.9107 and a macro-F1 of
0.8991. All handmade models exceeded 0.85 macro-F1, and the best CNN-BiLSTM
model reached the strongest overall result. This shows that the markerless pose
representation contains useful information for distinguishing correct from
false running technique when the running cycles are cleanly segmented.

\paragraph{RQ2 --- How much performance is lost when handmade cycles are
replaced by automatic autocorrelation-based detection?}
The loss is substantial. On the controlled fixed split, the reported autocorr
landmark baseline reached a test accuracy of 0.5455 and a macro-F1 of 0.5417,
compared with 0.9107 accuracy and 0.8991 macro-F1 for the handmade upper
bound. This corresponds to a drop of 0.3652 in accuracy and 0.3574 in
macro-F1. The fair 5-fold video-level CV check confirms that this loss is not
only a fixed-split artefact: under the same MLP classifier and threshold 0.5,
handmade cycles reached $0.860 \pm 0.016$ macro-F1, while autocorr cycles
reached $0.615 \pm 0.081$, leaving a fair gap of 0.245 macro-F1. A permutation
test on the autocorr MLP setting did not reach statistical significance
($p = 0.2475$), suggesting that autocorr-derived features are not yet robustly
separable under this simple classifier. Since the comparison keeps the general
landmark-representation family fixed while changing the cycle source, the loss
is attributed primarily to automatic cycle detection and boundary alignment.

\paragraph{RQ3 --- Which motion signal is most reliable for automatic cycle
detection?}
The answer depends on the criterion. In the signal ablation, \texttt{hip\_y}
achieved the highest downstream classification macro-F1 (0.7321) and accuracy
(0.7333). However, \texttt{contact\_diff} was retained as the main automatic
baseline signal because it is more directly tied to foot-contact dynamics and
supports foot-strike-to-foot-strike cycle interpretation. Thus, \texttt{hip\_y}
was the strongest signal by classification score in this ablation, while
\texttt{contact\_diff} was the most appropriate baseline signal for
semantically aligned cycle detection.

\paragraph{RQ4 --- Which architecture is most suitable for eight-phase pose
sequences?}
The answer is pipeline-dependent. On the handmade upper-bound pipeline, the
CNN-BiLSTM achieved the best macro-F1 (0.8991), indicating that local temporal
patterns and sequential modelling are useful when cycle alignment is reliable.
On the automatic autocorr pipeline, the CNN-BiLSTM was strongest at the default
0.5 threshold, but it became unstable under validation-tuned thresholding. The
MLP was therefore used as the more robust reported automatic baseline. This
suggests that higher model capacity is useful only when the cycle samples are
reliable; it does not compensate for noisy automatic segmentation.

\paragraph{RQ5 --- How do feature representation and frames-per-cycle affect
performance?}
The selected 13-landmark representation and the selected-landmark-plus-velocity
variant matched the best handmade macro-F1 (0.8991), while the full
33-landmark and lower-body-only variants remained close but did not improve the
overall macro-F1. This supports using a compact landmark representation as the
default. For temporal resolution, the eight-frame setting achieved the best
macro-F1 (0.8991), while 16, 24, and 32 frames did not yield a monotonic
improvement. The eight-phase representation is therefore both effective and
computationally economical for this dataset.

\paragraph{RQ6 --- How does the image-based representation compare with the
landmark-based representation under the autocorr pipeline?}
Under the controlled matched-cycle comparison, the landmark representation was
clearly more reliable. The autocorr landmark baseline reached a macro-F1 of
0.5417, while the image-based MobileNetV2--LSTM pipeline reached only 0.1000
macro-F1 and a ROC-AUC of 0.0592. The historical
\texttt{autocorr\_fixed} image-based milestone with validation accuracy 76.5\%
was not reproduced under the final fixed-split protocol. It is therefore treated
as development evidence and motivation for the organised image-based rerun, not
as a benchmark result.

\paragraph{RQ7 --- How do cycle aggregation, threshold selection, and pose/cycle
quality affect interpretation?}
Each factor materially affects interpretation. Aggregating cycle-level
predictions to the video level by majority vote improved the automatic landmark
baseline from 0.5139 to 0.5333 macro-F1, although the video-level test set was
small. Validation-tuned thresholds improved or preserved some clean handmade
results but overfit several automatic-pipeline models, especially the
CNN-BiLSTM. The quality analysis showed that higher heel visibility, lower
missing-ratio, and stronger autocorrelation were associated with better
classification accuracy. These findings confirm that evaluation should report
cycle-level and video-level behaviour, state whether thresholds are tuned on
validation, and account for pose and cycle quality.

% ---------------------------------------------------------------------
\section{Future Work}
\label{sec:concl-future}

The findings point to several future directions, ordered by their likely impact
on the automatic pipeline.

\begin{itemize}
    \item \textbf{Larger and more diverse dataset.} The validation and test sets
    are small, so future work should collect more in-the-wild videos across
    more runners, camera viewpoints, running speeds, clothing conditions,
    and recording environments.
    
    \item \textbf{Stronger foot-strike and cycle-boundary annotation.} Since the
    largest performance gap comes from replacing handmade cycles with
    automatic cycles, higher-quality foot-strike annotations should be
    collected to train and evaluate better boundary detectors.
    
    \item \textbf{Improved automatic cycle detection.} The current
    autocorrelation approach is interpretable but limited. Future work
    should compare it with learned boundary detectors, temporal
    segmentation models, and hybrid methods that combine periodicity,
    foot-contact cues, and confidence-based filtering.
    
    \item \textbf{Pose-quality-aware filtering and weighting.} The quality
    analysis indicates that errors concentrate in low-visibility or weakly
    periodic cycles. Future systems should explicitly use pose visibility,
    landmark missing-ratio, autocorrelation strength, and cycle-quality
    scores to reject, weight, or calibrate predictions.
    
    \item \textbf{Expert biomechanics validation.} The correct/false labels and
    recovered kinematic patterns should be checked against expert assessment
    so that the classification target is better grounded in running-form
    criteria.
    
    \item \textbf{More robust image-based modelling.} The controlled image rerun
    showed that RGB crops are unstable in the current implementation.
    Future work should improve runner tracking, bounding-box smoothing,
    crop stabilisation, and data augmentation before drawing conclusions
    about the value of image-based representations.
    
    \item \textbf{Pose--image fusion.} Rather than choosing between landmarks and
    crops, future models could fuse a compact landmark stream with a
    stabilised visual stream, combining the robustness of pose coordinates
    with visual cues such as silhouette and limb contour.
    
    \item \textbf{Video-level decision protocol.} Since aggregation can reduce
    cycle-level noise, future work should treat video-level prediction as a
    first-class evaluation target, compare majority vote with probability
    averaging and confidence voting, and report uncertainty when few cycles
    are available.
    
    \item \textbf{Deployment only after validation.} A real-time coaching or
    feedback tool is a possible long-term direction, but it should only be
    considered after the automatic pipeline is validated on larger,
    representative data and shown to be reliable across viewpoints and
    runner populations.
\end{itemize}

The final conclusion is therefore deliberately conservative. This thesis shows
that markerless pose data from in-the-wild running videos can support
running-technique classification when cycle segmentation is reliable. It also
shows that the current automatic system is limited mainly by upstream pose and
cycle quality. The next stage of work should therefore prioritise robust cycle
detection, larger validation data, and expert-labelled evaluation before moving
toward practical deployment.

% End of Chapter 7